\sectionguard\subsection{Augustine}\label{sec:fl-augustine}

Augustine of Hippo is the principal Western authority on the angels. He
wrote no treatise devoted to them, yet no Latin Father says so much about
them or says it in so ordered a way. The angels are the first work of God,
the light of the first day, the elder citizens of the one City of God;
some of them fell by their own will, and the place they left is filled
from redeemed men; the holy angels govern the visible creation under God,
bear his word to men, and refuse the sacrifice that belongs to him alone.
The teaching runs through his greatest works: the literal commentary on
Genesis (\emph{De Genesi ad litteram}), the \emph{City of God}
(\emph{De civitate Dei}), the third book \emph{On the Trinity}
(\emph{De Trinitate}), and the \emph{Enchiridion}, the handbook on faith,
hope, and charity written for Laurentius. Shorter works add the name of
their office (a sermon on Psalm~103), the limits of the demons' knowledge
(\emph{De divinatione daemonum}), and the rule of the honour paid to the
angels (\emph{De vera religione}).

Aquinas's treatise on the angels cites him in every question of \ST{I
qq.~50--64} except q.~53, and in every question of qq.~106--114
(\S\S\ref{sec:knowledge}, \ref{sec:creation}, \ref{sec:fall},
\ref{sec:government}). Augustine's own order runs from the creation of the
angels and their knowledge, through the division of the two societies and
the fall of the devil, to the repair of the heavenly city, the angels' part
in God's government, the worship they refuse and the worship they share,
and the vision they already enjoy.

\subsubsection{Spirit by nature, angel by office}\label{sec:aug-name}

Augustine teaches that the word \emph{angel} names what the angels do, not
what they are. Preaching on the verse ``Who makest thy angels spirits: and
thy ministers a burning fire'' (Ps 103[104]:4, Douay), he first sets down
the fact. Their appearing is hidden from our eyes, since they belong
\latin{in quadam republica magna imperatoris Dei}, in a great commonwealth
of God the emperor; \latin{tamen esse Angelos novimus ex fide, et multis
apparuisse scriptum legimus, et tenemus, nec inde dubitare fas nobis est}
--- yet we know by faith that the angels exist, we read that they have
appeared to many, and we hold it, nor may we doubt it. Then the
distinction:
\begin{quote}
\latin{Spiritus autem Angeli sunt; et cum spiritus sunt, non sunt angeli;
cum mittuntur, fiunt angeli. Angelus enim officii nomen est, non naturae.
Quaeris nomen huius naturae, spiritus est; quaeris officium, angelus est:
ex eo quod est, spiritus est; ex eo quod agit, angelus est.}
(\emph{Enarrationes in Psalmos} 103, s.~1.15)
\end{quote}
While they simply are spirits they are not yet angels; when they are sent,
they become angels. The Catechism of the Catholic Church opens its account
of who the angels are with this passage: ``\,`Angel' is the name of their
office, not of their nature. If you seek the name of their nature, it is
`spirit'; if you seek the name of their office, it is `angel': from what
they are, `spirit', from what they do, `angel'\,'' (CCC~329). Augustine
illustrates from human life: \latin{Nomen naturae homo, officii miles:
nomen naturae vir, officii praeco}. A man becomes a herald; a herald does
not become a man. So God takes those \latin{qui erant iam spiritus conditi
a creatore Deo}, spirits already made by the Creator, and makes them his
angels \latin{mittendo eos nuntiare quod iusserit}, by sending them to
announce what he has commanded. The burning fire of the verse he reads
first of the fire in the bush and of fire sent from heaven to do God's
bidding: \latin{cum esset, in natura sua erat; cum egit quod iussum est,
ministerium implevit} --- while it simply was, it was in its own nature;
when it did what it was bidden, it fulfilled a ministry. In the spiritual
sense God makes his preachers his angels, and Paul was sent to the carnal
Corinthians \latin{tamquam angelus de coelo ad terram}, as an angel from
heaven to earth; for \latin{fervens spiritu, ignis ardens est omnis
minister Dei}, every minister of God fervent in spirit is a burning fire
(s.~1.16).

The \emph{City of God} gives the same exegesis in the course of another
argument: ``He makes those who are by nature spirits His angels by
appointing them to the duty of bearing His messages. For the Greek word
angelos, which in Latin appears as `angelus,' means a messenger''
(\emph{De civitate Dei} XV.23). In the Latin, God makes them angels
\latin{iniungendo eis officium nuntiandi}, by enjoining on them the office
of announcing. The same passage leaves open whether the ``flaming fire''
names the bodies of the ministers or the love with which they ought to
burn, ``as with a spiritual fire'' --- a sign of the reserve
Augustine keeps on the question of angelic bodies (\S\ref{sec:aug-government}).

The literal commentary on Genesis adds why God has messengers at all. He
does not need them to learn about lower things, for he knows all things
simply and unchangeably; he has them for our sake and for theirs, because
to stand before him, to consult him about what lies below, and to obey his
commands \latin{bonum est eis in ordine propriae naturae atque
substantiae}, is good for them in the order of their own nature. By this
general name the whole heavenly city is called, \latin{quo nomine generali
universa illa superna civitas nuncupatur, quem primum diem conditum
existimamus} --- the city which Augustine takes to be the first day that
was made (\emph{De Genesi ad litteram} V.19.37). Aquinas keeps the same
derivation: \latin{Angelus nuntius dicitur}, and all the heavenly spirits
are called angels inasmuch as they manifest divine things (\ST{I q.~108
a.~5 ad~1}).

\subsubsection{The light of the first day}\label{sec:aug-creation}

The eleventh book of the \emph{City of God} begins the history of the two
cities with the angels, ``who constitute a large part of this city, and
indeed the more blessed part, since they have never been expatriated''
(XI.9). Genesis does not say plainly ``whether or when the angels were
created,'' yet they cannot have been omitted from the works from which God
rested; if they are mentioned, it is ``implicitly under the name of
`heaven' \dots\ or perhaps rather under the name of `light'.'' That they
were made, Scripture elsewhere says outright. The
canticle of the three children calls on ``all ye works of the Lord'' and
then names the angels among them (Dan 3:57--58); Psalm 148 bids the
angels and hosts praise God, ``for He commanded, and they were created''
(Ps 148:2, 5); ``Here the angels are most expressly and by divine
authority said to have been made by God.''

The argument for the day of their making proceeds by elimination. Job, in
Augustine's text, has God say, ``When the stars were made, the angels
praised me with a loud voice'' (Job 38:7). The angels therefore existed
before the stars, which were made on the fourth day, and the works of the
second and third days are named. ``There is no question, then, that
if the angels are included in the works of God during these six days, they
are that light which was called `Day,' and whose unity Scripture signalizes
by calling that day not the `first day,' but `one day'\,'' (Gen 1:5). The
second day and the rest ``are not other days; but the same `one' day is
repeated to complete the number six or seven.'' The doctrine follows from
the exegesis:
\begin{quote}
For when God said, ``Let there be light, and there was light,'' if we are
justified in understanding in this light the creation of the angels, then
certainly they were created partakers of the eternal light which is the
unchangeable Wisdom of God, by which all things were made, and whom we call
the only-begotten Son of God; so that they, being illumined by the Light
that created them, might themselves become light and be called ``Day,'' in
participation of that unchangeable Light and Day which is the Word of God,
by whom both themselves and all else were made. (XI.9)
\end{quote}
The Latin puts the whole in a clause: \latin{ut ea luce inluminati, qua
creati, fierent lux}, lit by the light by which they were made, so that
they might become light. ``The true Light, which lighteth every man that
cometh into the world'' (John 1:9) ``lighteth also every pure angel, that
he may be light not in himself, but in God; from whom if an angel turn
away, he becomes impure, as are all those who are called unclean
spirits.'' Their impurity is a privation, for \latin{Mali enim nulla
natura est; sed amissio boni mali nomen accepit} --- evil has no nature;
the loss of good has received the name of evil (XI.9). The next chapters
draw the consequence for the angels who stood. They ``were never at any
time or in any way darkness, but, as soon as they were made, were made
light; yet they were not so created in order that they might exist and
live in any way whatever, but were enlightened that they might live wisely
and blessedly'' (XI.11).

The commentary on Genesis gives the same reading in the terms of form and
formation. The light of the first day signifies the making of the
spiritual and intellectual creature, \latin{in qua natura intelleguntur
omnes Angeli sancti atque Virtutes}, in which nature all the holy Angels
and Virtues are understood. That is why Scripture does not add ``and God
made'' after ``Let there be light'': \latin{quia non primo cognovit
rationalis creatura conformationem suam, ac deinde formata est; sed in
ipsa sua conformatione cognovit, hoc est illustratione veritatis, ad quam
conversa formata est} --- the rational creature did not first know its
formation and then receive it; it knew it in the very act, by the
illumination of the truth to which it turned and so was formed
(\emph{De Gen.\ ad litt.}\ II.8.16). The spiritual creature, made in its
unformed beginning, was formed by turning to its Creator, \latin{a sua
quadam informitate ad Creatorem conversa atque formata} (IV.22.39).
Aquinas meets this reading as an objection when he determines that the
angels were created in grace, and answers that the lack of form preceded
their formation by nature and not by time (\ST{I q.~62 a.~3 obj.~1,
ad~1}; \S\ref{sec:q62}).

Augustine holds the creation of the angels with the world, not before it.
God \latin{creavit omnia simul}, as Ecclesiasticus says (Ecclus 18:1;
\emph{De Gen.\ ad litt.}\ IV.33.52), and among the things made at once he
set the spiritual creature above the bodily, because the spiritual can be
changed only through time and the bodily through time and place:
\latin{Spiritalem autem creaturam corporali praeposuit; quod spiritalis
tantummodo per tempora mutari posset, corporalis autem per tempora et
locos} (VIII.20.39). God, who moves neither in time nor place, moves the
created spirit through time and the body through time and place; Aquinas
answers with this text the objection that the angels, being above time,
were made from eternity (\ST{I q.~61 a.~2 ad~2}). Augustine knows another
opinion, which puts the angels before ``In the beginning'' and reads that
phrase as ``in the Word,'' the Beginning. He will not contend against it,
``chiefly because it gives me the liveliest satisfaction to find the
Trinity celebrated in the very beginning of the book of Genesis'':
\begin{quote}
Let each one, then, take it as he pleases; for it is so profound a
passage, that it may well suggest, for the exercise of the reader's tact,
many opinions, and none of them widely departing from the rule of faith.
At the same time, let none doubt that the holy angels in their heavenly
abodes are, though not, indeed, co-eternal with God, yet secure and
certain of eternal and true felicity. (XI.32)
\end{quote}
His own preference he states at the close of the book. The man of God
did not omit the angels, ``whether he included them in the words `in the
beginning,' because He made them first, or, which seems most likely,
because He made them in the only-begotten Word''; and ``heaven and earth''
name the whole creation, ``either as divided into spiritual and material,
which seems the more likely, or into the two great parts of the world''
(XI.33). Aquinas carries both opinions as Augustine does: the angels'
creation together with the corporeal world is ``the more probable'', and
the contrary ``is not to be deemed erroneous''; if it is held, ``In the
beginning'' must be read ``In the Son'' (\ST{I q.~61 a.~3};
\S\ref{sec:q61}). The Church has since defined what both opinions share,
that God \latin{simul ab initio temporis utramque de nihilo condidit
creaturam, spiritualem et corporalem, angelicam videlicet et mundanam}
(Lateran~IV, DS~800; \S\ref{sec:faith}).

One reading Augustine set aside. Some understood the angels by the waters
above the firmament (Gen 1:6--7), the waters below being visible water,
or the bad angels, or the nations; if that were so, ``then it does not
here appear when the angels were created, but when they were separated''
(XI.34). He had himself once written, in the last book of the
\emph{Confessions}, that the firmament was made between spiritual waters
above and bodily waters below, and later judged the phrase unconsidered:
\latin{non satis considerate dictum est; res autem in abdito est valde}
(\emph{Retractationes} II.6.2).

\subsubsection{Knowledge in the Word and in their own nature}\label{sec:aug-knowledge}

Augustine's doctrine of angelic knowledge grows out of the formulas of
Genesis~1. Each work is announced three times: \latin{Et dixit Deus:
Fiat}; \latin{Et sic est factum}; \latin{Fecit Deus}. The first refers the
work to the eternity of the Word; the second to the knowledge of it made in the
intellectual creature; the third to the thing made in its own kind. So the
creation of the heaven was first in the Word of God, as in begotten
Wisdom; \latin{deinde facta est in creatura spiritali, hoc est in
cognitione Angelorum secundum creatam in illis sapientiam}, then it was
made in the spiritual creature, that is, in the knowledge of the angels,
according to the wisdom created in them; and then the heaven itself was
made in its own kind (\emph{De Gen.\ ad litt.}\ II.8.16, 19). The angels'
knowledge of every other work is therefore prior to the work, since they
were made first: as the reason by which a creature is made is prior in
the Word to the creature made, so the knowledge of that reason comes first
in the intellectual creature \latin{quae peccato tenebrata non est}, not
darkened by sin, and then the making of the creature follows (II.8.17). The angels do not see sensible things as cattle do,
\latin{solo sensu corporis}; if they use any such sense, they recognize
by it what they know better within, in the Word by whom they are lit. Nor
did they advance to wisdom as we do, from the visible to the invisible
(Rom 1:20), \latin{qui ex quo creati sunt, ipsa Verbi aeternitate sancta
et pia contemplatione perfruuntur} --- from the moment they were made they
enjoy the eternity of the Word in holy and devout contemplation, and
looking thence upon creatures they approve what is rightly done and
reprove sin (II.8.17). God showed his holy angels beforehand what he was
going to make, and they knew his mind only so far as he showed it: ``For
who hath known the mind of the Lord?'' (Rom 11:34; II.8.18). Aquinas
answers with this chapter the objection to his determination that the
angels do not understand by species drawn from things (\ST{I q.~55 a.~2
ad~1}), makes it the \emph{sed contra} of his article on whether their
intellect is ever in potency (q.~58 a.~1), and takes from it the twofold
procession of things from the Word, into the angelic mind and into their
own natures (q.~56 a.~2).

From here comes the figure for which Augustine is best known in
angelology, the day, the evening, and the morning of the angelic mind. If
the light of the first day is spiritual, it was made after darkness as
the spiritual creature turned from its unformed state to the Creator,
\latin{ita et post vesperam fiat mane}, so after evening morning is made
when, having known its own nature, \latin{qua non est quod Deus}, by which it is not what God is, it
refers itself to the praise of the light which God is (IV.22.39). Evening
is the angels' knowledge of a creature in itself; morning is their turning
from that knowledge to praise, in which they receive from the Word the
knowledge of the next work to be made. The great difference between
knowing a thing in the Word and knowing it in its own nature fixes the
names, the one belonging to the day and the other to the evening
(IV.23.40; \S\ref{sec:q58}). Beside the light seen in the Word, every knowledge of a creature in
itself may be called night; beside the error of those who do not know
even the creature, it may be called day. So the life of the faithful is
light beside the life of unbelief (Eph 5:8), and night beside the day in
which, \latin{aequales Angelis facti}, made equal to the angels, we shall
see God as he is (2 Pet 1:19; IV.23.40). The holy angels, \latin{in
quibus prima omnium creata est sapientia}, in whom created wisdom was the
first of all things made, always see the face of God and enjoy his only
Word; they therefore know the whole creation first in the Word, where the eternal
reasons of all things are, and then in the creature itself, \latin{quam
sic noverunt tamquam infra despicientes}, which they know as looking down
upon it, and which they refer to the praise of him in whose truth they see
its reasons. At once it becomes morning, because the angelic knowledge
does not stay in the thing made. Had the angelic nature turned instead to
itself and taken more delight in itself than in him by whose participation
it is blessed, \latin{intumescens superbia caderet, sicut diabolus} ---
swelling with pride, it would fall, as the devil fell (IV.24.41). The
doctrine of the angels' knowledge and the doctrine of the fall are one
doctrine seen from two sides.

For this reason the six days have no night. \latin{Tunc enim nox ad diem
pertinet, non dies ad noctem}: the night belongs to the day and not the
day to the night, when the high and holy angels refer what they know of
the creature to the honour and love of him in whom they contemplate its
eternal reasons; \latin{eaque concordissima contemplatione sunt unus
dies, quem fecit Dominus, cui coniungetur Ecclesia de hac peregrinatione
liberata} --- by that most concordant contemplation they are the one day
which the Lord has made, to which the Church, freed from this
pilgrimage, will be joined (IV.25.42; Ps 117[118]:24). The day is
repeated through the works \latin{non circuitu corporali, sed cognitione
spiritali}, not by a bodily circuit but by spiritual knowledge (IV.26.43). Augustine will not have this taken as a figure laid over the
real days. Christ is called light properly and stone figuratively; so
here, \latin{Ubi enim melior et certior lux, ibi verior etiam dies: cur
ergo non et verior vespera et verius mane?} --- where the light is better
and more certain, there the day is truer; why not then a truer evening and
a truer morning? (IV.28.45).

Evening, morning, and day are not successive in the angels as they are on
earth. The angels of the high heavens hold all three together
(IV.29.46), and in the heavenly homeland all three are always present:
\latin{ubi semper est dies in contemplatione incommutabilis veritatis,
semper vespera in cognitione in seipsa creaturae, semper mane etiam ex
hac cognitione in laude Creatoris} --- there it is always day in the
contemplation of unchangeable truth, always evening in the knowledge of the
creature in itself, always morning in the praise of the Creator that
rises from that knowledge (IV.30.47). While the works were being made,
the order was of causes and not of delays: \latin{Mens vero angelica pura
caritate inhaerens Verbo Dei}, the angelic mind cleaving to the Word of
God by pure charity, saw the works in the Word before they were made
(IV.32.49). Of the first light Augustine says,
\latin{adhaesitque creanti luci lux creata, videns illam et se in illa,
id est rationem qua facta est. Vidit etiam se in se} --- the created
light clung to the creating light, seeing it, and itself in it, that is,
the reason by which it was made; it saw itself also in itself (IV.32.50).
The six days are ordered \latin{prius et posterius habens in connexione
creaturarum, in efficacia vero Creatoris omnia simul}, with before and
after in the connection of creatures, and all at once in the efficacy of
the Creator; and God's rest on the seventh day is his rest \latin{in
conspectu Angelorum suorum}, in the sight of his angels (IV.35.56).

The \emph{City of God} gives the doctrine to a wider audience (XI.7,
29). The holy angels know God
\begin{quote}
not by audible words, but by the presence to their souls of immutable
truth, i.e., of the only-begotten Word of God; and they know this Word
Himself, and the Father, and their Holy Spirit, and that this Trinity is
indivisible, and that the three persons of it are one substance, and that
there are not three Gods but one God; and this they so know that it is
better understood by them than we are by ourselves. (XI.29)
\end{quote}
They know the creature in the wisdom of God, ``as if in the art by which
it was created,'' and ``know themselves better in God than in themselves''
(XI.29). The English ``noonday knowledge'' renders
Augustine's \latin{tamquam in diurna cognitione}, a day-knowledge, set
against \latin{in se ipsis autem tamquam in uespertina}, an evening
knowledge in themselves.

The angels' knowledge is not confined to creation. They knew from the ages
the mystery of the kingdom of heaven revealed at the appointed time for our
salvation, since the seed that came in its time was ``ordained by angels in
the hand of a mediator'' (Gal 3:19); and the Apostle's word that the
manifold wisdom of God is made known ``to the principalities and powers in
heavenly places through the church'' (Eph 3:10) means that the Church was
first there, where the Church of the resurrection will be gathered
(\emph{De Gen.\ ad litt.}\ V.19.38--39; 1 Tim 3:16). Aquinas cites
this text in the objection of his article on the angels' knowledge of the
mysteries of grace and again of the demons' knowledge (\ST{I q.~57 a.~5
obj.~1}; q.~64 a.~1 obj.~4).

Between this knowledge and the demons' Augustine draws the line that
Aquinas follows. The good angels ``have a more certain knowledge even of those temporal and mutable
things, because they contemplate their principles and causes in the word
of God''; the demons ``only foresee a larger part of the future than men
do, by reason of their greater acquaintance with the signs which are
hidden from us. Sometimes, too, it is their own intentions they predict.
And, finally, the demons are frequently, the angels never, deceived''
(\emph{De civ.\ Dei} IX.22). Aquinas receives Augustine at each step:
in the division of morning and evening knowledge (\ST{I q.~58 aa.~6--7}),
in the angels' conjecture of secret thoughts (q.~57 a.~4), and in the
knowledge left to the demons (q.~64 a.~1; \S\S\ref{sec:q58},
\ref{sec:q64}).

\subsubsection{The two societies of angels}\label{sec:aug-societies}

When God ``divided the light from the darkness'' (Gen 1:4), Augustine sees
the separation of the holy angels from the unclean. ``To me it does not
seem incongruous with the working of God, if we understand that the angels
were created when that first light was made, and that a separation was
made between the holy and the unclean angels, when, as is said, `God
divided the light from the darkness; and God called the light Day, and
the darkness He called Night.' For He alone could make this
discrimination, who was able also before they fell, to foreknow that they
would fall, and that, being deprived of the light of truth, they would
abide in the darkness of pride'' (\emph{De civ.\ Dei} XI.19). He reads
the wording closely. After ``Let there be light'' comes ``God saw the
light that it was good''; after the division no such approval follows,
``lest both should be designated good, while one of them was evil, not by
nature, but by its own fault.'' So ``the light alone received the
approbation of the Creator, while the angelic darkness, though it had
been ordained, was yet not approved'' (XI.20). Aquinas gives the answer of
this chapter to the objection that the darkness of Genesis~1 signifies
demons wicked from their creation: the division is made by God's
foreknowledge (\ST{I q.~63 a.~5 ad~2}).

The same exegesis returns at the close of the book, where the fall is
taught from the Second Epistle of Peter: ``God spared not the angels that
sinned, but delivered them, drawn down by infernal ropes to the lower
hell, unto torments, to be reserved unto judgment'' (2 Pet 2:4, Douay).
The holy angels are rightly called light, for even we who live by faith
are light in the Lord (Eph 5:8), and the apostates are darkness. Then
Augustine sets the two societies against each other in a run of
antitheses:
\begin{quote}
we understand these two societies of angels,---the one enjoying God, the
other swelling with pride; the one to whom it is said, ``Praise ye Him,
all His angels,'' the other whose prince says, ``All these things will I
give Thee if Thou wilt fall down and worship me;'' the one blazing with
the holy love of God, the other reeking with the unclean lust of
self-advancement. \dots\ the one, at God's pleasure,
tenderly succoring, justly avenging,---the other, set on by its own pride,
boiling with the lust of subduing and hurting; the one the minister of
God's goodness to the utmost of their good pleasure, the other held in by
God's power from doing the harm it would; the former laughing at the
latter when it does good unwillingly by its persecutions, the latter
envying the former when it gathers in its pilgrims. (XI.33)
\end{quote}
The two communities are ``the one both by nature good and by will
upright, the other also good by nature but by will depraved'' (XI.33).
Aquinas cites this chapter in the objection that the demons can
know no truth, since the good angels are separated from them as light
from darkness, and answers it with the three ways in which the demons
still know truth: by the subtlety of their nature, by revelation from the
holy angels, and by long experience (\ST{I q.~64 a.~1 obj.~5, ad~5}).

The division leaves the order of natures standing and reorders merit. By
nature the angels
stand above men, ``the immortal such as the angels, above the mortal,
such as men''; ``but of such consequence in rational natures is the
weight, so to speak, of will and of love, that though in the order of
nature angels rank above men, yet, by the scale of justice, good men are
of greater value than bad angels'' (XI.16). Nor are there four cities,
two of angels and two of men, ``but rather two in all, one composed of
the good, the other of the wicked, angels or men indifferently'' (XII.1).
Aquinas cites this sentence for the one hierarchy of all rational
creatures under God (\ST{I q.~108 a.~1}; \S\ref{sec:q108}).

The twelfth book asks where the two wills came from. The difference arose
``not from a difference in their nature and origin \dots\ but from a
difference in their wills and desires'':
\begin{quote}
While some steadfastly continued in that which was the common good of
all, namely, in God Himself, and in His eternity, truth, and love; others,
being enamored rather of their own power, as if they could be their own
good, lapsed to this private good of their own, from that higher and
beatific good which was common to all \dots\ they became proud,
deceived, envious. The cause, therefore, of the blessedness of the
good is adherence to God. (XII.1)
\end{quote}
The fault proves the nature: ``when we say that it is a fault of the
angelic creature that it does not cleave to God, we hereby most plainly
declare that it pertained to its nature to cleave to God'' (XII.1). The
commentary on Genesis names the two loves that divided the angels before
they divided men, \latin{alter sanctus \dots\ alter immundus; alter
socialis, alter privatus}, the one holy and social, the other unclean and
private, the one subject to God and the other his rival, and says of them
that \latin{praecesserunt in Angelis; alter in bonis, alter in malis},
they went before in the angels, one in the good, the other in the evil,
and then distinguished the two cities founded in mankind (\emph{De Gen.\
ad litt.}\ XI.15.20).

The cause of the evil angels' misery is that they turned from him who
supremely is to themselves, ``and this vice, what else is it called than
pride? For `pride is the beginning of sin''' (XII.6; Ecclus 10:15). Of
the evil will itself there is no efficient cause. ``If the further
question be asked, What was the efficient cause of their evil will? there
is none'' (XII.6). No other will can make it evil, and no lower thing,
for every nature is good, and ``when the will
abandons what is above itself, and turns to what is lower, it becomes
evil---not because that is evil to which it turns, but because the
turning itself is wicked'' (XII.6). Augustine concludes:
\begin{quote}
Let no one, therefore, look for an efficient cause of the evil will; for
it is not efficient, but deficient, as the will itself is not an effecting
of something, but a defect. For defection from that which supremely is,
to that which has less of being,---this is to begin to have an evil will.
Now, to seek to discover the causes of these defections,---causes, as I
have said, not efficient, but deficient,---is as if some one sought to
see darkness, or hear silence. (XII.7)
\end{quote}
In Latin the paradox is a play of two words, \latin{non enim est efficiens
sed deficiens, quia nec illa effectio sed defectio}. The will could
not become evil unwilling, and its defections are justly punished,
``being not necessary, but voluntary'' (XII.8). Pride in particular ``is not the fault of him who delegates power, nor of power
itself, but of the soul that is inordinately enamored of its own power,
and despises the more just dominion of a higher authority'' (XII.8). This
is the analysis Aquinas gives of the angel's first sin, a sin in a
spiritual good sought ``not according to proper measure or rule'' (\ST{I
q.~63 a.~1 ad~4}; \S\ref{sec:q63}).

The good will has a cause, and its cause is God. Augustine argues by
dilemma. If the good angels made their own will good, they did so either
without will, and it was not their doing, or with a will already good,
and then they had it already:
\begin{quote}
And who made this will, which already they had, but He who created them
with a good will, or with that chaste love by which they cleaved to Him,
in one and the same act creating their nature, and endowing it with
grace? And thus we are driven to believe that the holy angels never
existed without a good will or the love of God. (XII.9)
\end{quote}
The Latin has \latin{simul eis et condens naturam et largiens gratiam},
at once establishing their nature and bestowing grace. Of the angels who fell Augustine gives two
possibilities and chooses neither: ``These angels, therefore, either
received less of the grace of the divine love than those who persevered
in the same; or if both were created equally good, then, while the one
fell by their evil will, the others were more abundantly assisted, and
attained to that pitch of blessedness at which they became certain they
should never fall from it'' --- in his Latin, \latin{aut minorem
acceperunt diuini amoris gratiam \dots\ illi amplius adiuti}. He then
applies to the angels the verse on charity:
\begin{quote}
We must therefore acknowledge, with the praise due to the Creator, that
not only of holy men, but also of the holy angels, it can be said that
``the love of God is shed abroad in their hearts by the Holy Ghost, which
is given unto them.'' And that not only of men, but primarily and
principally of angels it is true, as it is written, ``It is good to draw
near to God.'' And those who have this good in common, have, both with Him
to whom they draw near, and with one another, a holy fellowship, and form
one city of God---His living sacrifice, and His living temple. (XII.9;
Rom 5:5; Ps 72[73]:28)
\end{quote}
Aquinas makes the question and answer of this chapter the \emph{sed
contra} of his article that the angels were created in grace, which he
calls ``more probable, and more in keeping with the sayings of holy men''
(\ST{I q.~62 a.~3}; \S\ref{sec:q62}); he cites the same chapter for
charity poured into the angels by the Holy Spirit (q.~60 a.~5 obj.~4;
q.~108 a.~4 obj.~2), and for the one society of angels and men into which
the saints are taken (q.~108 a.~8).

\subsubsection{The fall of the devil}\label{sec:aug-fall}

Augustine's teaching on the devil's fall is built on one verse of the
Gospel: ``He was a murderer from the beginning: and he stood not in the
truth, because truth is not in him'' (John 8:44, Douay). He reads it
against the Manichees, who gave the devil a nature of his own from an
evil principle. ``The Lord did not say, `The devil was naturally a
stranger to the truth,' but `The devil abode not in the truth,' by which
He meant us to understand that he had fallen from the truth, in which, if
he had abode, he would have become a partaker of it, and have remained in
blessedness along with the holy angels'' (\emph{De civ.\ Dei} XI.13);
his not abiding in the truth is the cause of its not being in him, and not
the reverse (XI.14). ``The devil sinneth from the beginning'' (1 John 3:8) cannot
mean that he was made sinful, ``for if sin be natural, it is not sin at
all.'' The prophets speak against it: the king of Babylon, ``How art thou
fallen, O Lucifer, son of the morning!'' (Isa 14:12), and the prince of
Tyre, ``Thou hast been in Eden, the garden of God'' and ``Thou wast
perfect in thy ways'' (Ezek 28:13, 15), in whose persons the devil is
addressed. ``And if these passages cannot well be otherwise interpreted,
we must understand by this one also, `He abode not in the truth,' that he
was once in the truth, but did not remain in it'' (XI.15). He sinned
``from the beginning of his sin, when by his pride he had once commenced
to sin.'' The dragon ``made to be a sport to His angels'' (Job 40:14;
Ps 103[104]:26) was doomed to that mockery after his sin, and ``how much
more, then, is this angelic nature,
which surpasses in dignity all else that He has made, the handiwork of the
Most High!'' (XI.15).

The Church confesses the same truth in the words of the Fourth Lateran
Council: the devil and the other demons were created good by nature and
became evil of themselves (DS~800; \S\ref{sec:faith}). Augustine draws
from it the lesson of providence. God, ``as He is the supremely good
Creator of good natures, so is He of evil wills the most just Ruler; so
that, while they make an ill use of good natures, He makes a good use even
of evil wills. Accordingly, He caused the devil (good by God's creation,
wicked by his own will) to be cast down from his high position, and to
become the mockery of His angels,---that is, He caused his temptations to
benefit those whom he wishes to injure by them'' (XI.17). The devil's sin
is spiritual: ``though
we cannot call the devil a fornicator or drunkard, or ascribe to him any
sensual indulgence (though he is the secret instigator and prompter of
those who sin in these ways), yet he is exceedingly proud and envious''
(XIV.3); Aquinas makes this the \emph{sed contra} of his article that only
pride and envy can be in the angels (\ST{I q.~63 a.~2}).

The eleventh book of the commentary on Genesis gives the fullest
treatment. The devil's nature is good, and it still lives; his evil is a
will moved inordinately, preferring lower goods to higher, so that a
rational spirit, delighting in its own power for its excellence,
\latin{tumesceret superbia, per quam caderet a beatitudine spiritalis
paradisi, et invidentia contabesceret} --- swelled with pride, by which it
fell from the beatitude of the spiritual paradise, and wasted away with
envy (\emph{De Gen.\ ad litt.}\ XI.13.17).
Some said that the devil fell because he envied man made in God's image.
Augustine answers that envy follows pride and does not precede it, and
defines both:
\begin{quote}
\latin{Cum igitur superbia sit amor excellentiae propriae, invidentia
vero sit odium felicitatis alienae, quid unde nascatur satis in promptu
est. \dots\ Superbiendo igitur invidus, non invidendo quisque superbus
est.} (XI.14.18)
\end{quote}
Pride is the love of one's own excellence, envy the hatred of another's
happiness; one becomes envious by being proud, and not proud by being
envious. Pride is the beginning of all sin, and with it the Apostle's
``root of all evils,'' avarice taken in its general sense as the desire
for more than is due for the sake of one's own excellence: \latin{Hac enim
et diabolus cecidit, qui utique non amavit pecuniam, sed propriam
potestatem} --- by this the devil fell, who did not love money but his
own power (XI.15.19; 1 Tim 6:10). Aquinas's order of the angel's sins,
pride first and envy after, is Augustine's (\ST{I q.~63 a.~2};
\S\ref{sec:q63}).

When he fell, Scripture does not say. That he fell before he envied man
is plain, and Augustine judges it reasonable to think that he fell from
the very beginning of time:
\begin{quote}
\latin{Non autem frustra putari potest, ab initio temporis diabolum
superbia cecidisse, nec fuisse ullum antea tempus quo cum Angelis sanctis
pacatus vixerit et beatus, sed ab ipso exordio creaturae a suo creatore
apostatasse.} (XI.16.21)
\end{quote}
On this reading both clauses of John 8:44 speak of a beginning: he was a
murderer from the beginning of man, since he killed the first man, and he
did not stand in the truth from the beginning of his own creation,
\latin{qui staret si stare voluisset}, who would have stood had he willed
to stand. The reason is the certainty of the blessed. The holy angels are
not uncertain of their eternal life, \latin{Nam quomodo beati, si
incerti?}; how then could the devil have lived blessed among them without
foreknowing his sin and punishment? (XI.17.22). Some answered that the
devil was not of the supercelestial angels but of those made lower in the
world and distributed through its offices, and that such angels might be
blessed in a lower way, by hope. Augustine reports the opinion without
adopting it: \latin{mihi autem, fateor, unde hoc secundum Scripturas
asseram, non interim occurrit} --- for my part, I confess, I do not yet
see how I could assert this from the Scriptures (XI.19.26).

Created evil the devil was not. Those who read Job's ``beginning of the
creation of God'' so contradict ``God saw all the things that he had
made, and they were very good'' (Gen 1:31) and the justice of the
everlasting fire ``prepared for the devil and his angels'' (Matt 25:41);
what is punished in him is \latin{nullo modo \dots\ naturam quam Deus
creavit, sed malam propriam voluntatem}, not at all the nature God made,
but his own evil will (XI.21.28). Then comes Augustine's most exact
sentence on the fall:
\begin{quote}
\latin{Sed factus continuo se a luce veritatis avertit, superbia tumidus,
et propriae potestatis delectatione corruptus: unde beatae atque angelicae
vitae dulcedinem non gustavit, quam non utique acceptam fastidivit, sed
nolendo accipere deseruit, et amisit.} (XI.23.30)
\end{quote}
Made, he turned at once from the light of truth, swollen with pride and
corrupted by delight in his own power; he never tasted the sweetness of
the blessed and angelic life, which he did not receive and then despise,
but deserted and lost by refusing to receive it. He fell \latin{non ex eo
quod acceperat \dots\ sed ex eo quod acciperet, si subdi voluisset Deo},
not from what he had received, but from what he would have received had he
been willing to be subject to God. The oracles of Isaiah and Ezekiel, he
adds, fit the devil's body, the multitude of the wicked of which he is the
head, better than the head himself (XI.24.31--25.32).

He closes the inquiry with a \latin{complexio} of four possibilities:
the devil fell by impious pride from the beginning of his creation, from
the beatitude he would have received had he willed; or there are angels
of a lower ministry in this world among whom he lived, and from whom he
fell \latin{tamquam archangelus}, as an archangel, with the angels subject
to him, if that can be asserted at all; or the holy angels, if the devil
once lived among them, received the certainty of their perseverance only
after his fall; or else one must ask by what desert he and his companions
were set apart before their sin, so that they did not foreknow their fall
while the others were certain of their perseverance. Whichever is true, two
things are certain: the sinning angels are thrust into the dark prison of
this air to be reserved for judgment (2 Pet 2:4), and the holy angels'
eternal life is not uncertain, nor will ours be, when after the
resurrection we are joined to them (XI.26.33). The \emph{Enchiridion}
states the outcome in the language of the Creed. The rebel angels were
not all propagated from one, as men are from Adam, and so were not bound
by an inherited guilt; ``the rest remained steadfast in piety and
obedience to their Lord, and obtained, what before they had not enjoyed, a
sure and certain knowledge of their eternal safety, and freedom from the
possibility of falling'' (\emph{Enchiridion} 28); ``what catholic
Christian does not know that no new devil will ever arise among the good
angels, as he knows that this present devil will never again return into
the fellowship of the good?'' (\emph{De civ.\ Dei} XI.13). The Catechism teaches the same irrevocability:
``There is no repentance for the angels after their fall'' (CCC~393).

Augustine and Aquinas answer the question of the moment differently.
Augustine leaves the four possibilities open and inclines, in the
commentary on Genesis, to a fall at the very beginning of the devil's
creation, though never to a creation in evil; in the \emph{City of God}
he reads Isaiah and Ezekiel to say that the devil was once in the truth
and did not remain (XI.15). Aquinas records that Augustine does not count
the opinion of a devil wicked from the first instant of his creation
Manichaean, but judges that it ``was reasonably rejected by the masters as
erroneous''; he takes the sentence of \emph{City of
God} XI.15 as the \emph{sed contra} of his article on the interval and
determines as ``more probable'' that the devil sinned at once after the
first instant of his creation (\ST{I q.~63 aa.~5--6}; \S\ref{sec:q63};
Appendix~\ref{app:first-sin}). The opinion Augustine
reports, that the devil and his angels were of the lower orders set over
the world, reappears in Aquinas as Damascene's, against Gregory's that the
first of them was the highest of all (\ST{I q.~63 a.~7}; q.~61 a.~4
ad~2).

The devil fell by his own will, and he tempts only by God's leave. He
was not permitted to tempt the woman except through the serpent, nor the
man except through the woman, and his power of temptation is bounded:
\latin{tentandi voluntatem habet diabolus, in potestate autem nec ut
faciat habet, nec quo modo faciat} --- the devil has the will to tempt,
but neither whether he does it nor how is in his power (\emph{De Gen.\ ad
litt.}\ XI.27.34). God does this through the holy angels, \latin{subdens
videlicet angelos malos Angelis bonis, ut malorum improbitas, non quantum
nititur, sed quantum sinitur possit} --- subjecting the evil angels to the
good, so that the wickedness of the evil can do not as much as it strives
to do but as much as it is allowed (XI.22.29); and it is so until the
saints judge angels (1 Cor 6:3).

\subsubsection{The restoration of the heavenly city}\label{sec:aug-restoration}

The angels who fell are not restored; the part of the rational creation
that fell among men is restored in part, and fills the place the angels
left. The \emph{Enchiridion} gives the plan of God:
\begin{quote}
And so it pleased God, the Creator and Governor of the universe, that,
since the whole body of the angels had not fallen into rebellion, the
part of them which had fallen should remain in perdition eternally \dots\
but that, on the other hand, mankind, who constituted the
remainder of the intelligent creation, having perished without exception
under sin, both original and actual, and the consequent punishments,
should be in part restored, and should fill up the gap which the
rebellion and fall of the devils had left in the company of the angels.
(\emph{Enchiridion} 29)
\end{quote}
The promise that the saints ``are equal to the angels'' (Luke 20:36)
carries it, and ``the Jerusalem which is above, which is the mother of us
all, the city of God, shall not be spoiled of any of the number of her
citizens, shall perhaps reign over even a more abundant population''
(\emph{Enchiridion} 29). The number of the saints and of the demons is
unknown to us and present to the Creator, who ``ordereth all things in
measure, and number, and weight'' (Wis 11:21). The last book of the
\emph{City of God} states the same plan: God gathers from the condemned
race \latin{tantum populum gratia sua \dots\ ut inde suppleat et instauret
partem, quae lapsa est angelorum}, so great a people by his grace as to
supply and restore the fallen part of the angels, so that the beloved
city ``is not defrauded of the full number of its citizens, but perhaps
may even rejoice in a still more overflowing population'' (\emph{De civ.\
Dei} XXII.1). Aquinas assumes the same when he reasons that, if the chief
of the devils was of the highest order, some fell from every order, ``just
as men are taken up into every order to supply for the angelic ruin''
(\ST{I q.~63 a.~9 ad~3}); and his article that men are taken up into the
angelic orders rests on the one society of angels and men that Augustine
teaches (q.~108 a.~8; \S\ref{sec:q108}).

The \emph{Enchiridion} places all this within the Creed. After the Holy
Spirit the Creed names the holy Church, and the Church must be understood
whole:
\begin{quote}
not that part of it only which wanders as a stranger on the earth \dots\
but that part also which has always from its creation remained steadfast
to God in heaven, and has never experienced the misery consequent upon a
fall. This part is made up of the holy angels, who enjoy uninterrupted
happiness; and (as it is bound to do) it renders assistance to the part
which is still wandering among strangers: for these two parts shall be
one in the fellowship of eternity, and now they are one in the bonds of
love, the whole having been ordained for the worship of the one God.
(\emph{Enchiridion} 56)
\end{quote}
Neither part of the Church wishes to be worshipped instead of God, and
the order of the Creed itself proves the divinity of the Holy Spirit: if
he were a creature he would belong to the Church, ``to that part of it
which is in the heavens,'' and would be part of a temple rather than the
God who has one (56). Christ did not die for the angels. ``Yet what was done
for the redemption of man through His death was in a sense done for the
angels, because the enmity which sin had put between men and the holy
angels is removed, and friendship is restored between them, and by the
redemption of man the gaps which the great apostasy left in the angelic
host are filled up'' (61). The holy angels, ``taught by God, in the
eternal contemplation of whose truth their happiness consists, know how
great a number of the human race are to supplement their ranks''; and so
``all things are gathered together in one in Christ, both which are in
heaven and which are on earth'' (62; Eph 1:10; Col 1:20).

On the ranks of that society Augustine confesses ignorance. The whole
company is called by the general name of angels, for ``to which of the
angels said God at any time, Sit on my right hand?'' (Heb 1:13); yet some
are called archangels; whether the archangels are the ``hosts'' of the
psalm (Ps 148:2), and what the four names of the Apostle signify,
``thrones, or dominations, or principalities, or powers'' (Col 1:16),
\begin{quote}
let those who are able answer these questions, if they can also prove
their answers to be true; but as for me, I confess my ignorance. I am not
even certain upon this point: whether the sun, and the moon, and all the
stars, do not form part of this same society, though many consider them
merely luminous bodies, without either sensation or intelligence.
(\emph{Enchiridion} 58)
\end{quote}
The Latin is briefer and stronger: \latin{dicant qui possunt, si tamen
possunt probare quod dicunt: Ego me ista ignorare confiteor}; and the
second name of the psalm is \latin{Virtutes}. Augustine does not rank the
orders. He knows that there are degrees in the heavenly city, for he
teaches that in the city to come ``no inferior shall envy any superior,
as now the archangels are not envied by the angels'' (\emph{De civ.\ Dei}
XXII.30); but the ordering of the nine choirs comes to the West from
Dionysius and Gregory, and Aquinas determines it after them (\ST{I q.~108
aa.~5--6}; \S\S\ref{sec:ch}, \ref{sec:q108}; Appendix~\ref{app:orders}).
The same reserve covers the bodies in which angels appeared: how they
were seen and touched, so that the patriarchs washed their feet and Jacob
wrestled with one, and how they speak within the soul, ``the angel that
spoke in me'' (Zach 1:9), is a question that is ``a useful exercise for
the intellect, if the discussion be kept within proper bounds, and if we
avoid the error of supposing ourselves to know what we do not know''
(\emph{Enchiridion} 59). More necessary is to discern Satan when he
transforms himself into an angel of light (2 Cor 11:14; 60).

\subsubsection{Angels and demons in the government of the world}\label{sec:aug-government}

The third book \emph{On the Trinity} asks how God appeared to the
fathers. Before answering, Augustine lays down how God governs the bodily
world at all, and turns to the heavenly country:
\begin{quote}
For there the will of God, ``who maketh His angels spirits, and His
ministers a flaming fire,'' presiding among spirits which are joined in
perfect peace and friendship, and combined in one will by a kind of
spiritual fire of charity \dots\ thence diffuses itself through all
things by certain most perfectly ordered movements of the creature; first
spiritual, then corporeal \dots\ But as the more gross and
inferior bodies are governed in due order by the more subtle and powerful
ones, so all bodies are governed by the living spirit; and the living
spirit devoid of reason, by the reasonable living spirit; and the
reasonable living spirit that makes default and sins, by the living and
reasonable spirit that is pious and just; and that by God Himself, and so
the universal creature by its Creator. (\emph{De Trinitate} III.4.9)
\end{quote}
``And so it comes to pass that the will of God is the first and the
highest cause of all corporeal appearances and motions. For nothing is
done visibly or sensibly, unless either by command or permission from the
interior palace, invisible and intelligible, of the supreme Governor''
--- \latin{de interiore invisibili atque intellegibili aula summi
Imperatoris} --- ``in that far-reaching and boundless commonwealth of the
whole creature'' (III.4.9). This is the text by which the \emph{Summa}
teaches that all bodies are ruled by the angels: Aquinas makes it the
\emph{sed contra} of the article that the corporeal creature is governed
by the angels (\ST{I q.~110 a.~1}), and cites it again for the ordering
of the orders and for the government of the demons by the good angels
(q.~108 a.~6; q.~109 a.~4 \emph{sed contra}; \S\ref{sec:q110}). The
commentary on Genesis states
the same rule: \latin{sublimibus Angelis Deo subdite fruentibus, et Deo
beate servientibus, subdita est omnis natura corporea, omnis irrationalis
vita, omnis voluntas vel infirma vel prava} --- to the high angels who
enjoy God in subjection to him and serve him in blessedness, all bodily
nature is subject, all irrational life, every will weak or depraved;
they contemplate God without place or time and carry out his commands in
lower things (\emph{De Gen.\ ad litt.}\ VIII.24.45). God thus presides over his whole
creation by a twofold work of providence, \latin{naturis, ut fiant;
voluntatibus autem, ut sine suo iussu vel permissu nihil faciant}, over
natures that they may be, over wills that they may do nothing without his
command or permission (VIII.24.45). In the
book of eighty-three questions Augustine puts the rule in a sentence that
Aquinas quotes: \latin{unaquaeque res visibilis in hoc mundo habet
potestatem angelicam sibi praepositam, sicut aliquot locis divina
Scriptura testatur} --- every visible thing in this world has an angelic
power set over it, as divine Scripture testifies in several places
(\emph{De diversis quaestionibus octoginta tribus} q.~79.1; \ST{I q.~110
a.~1}).

Miracles belong to the same government. The same power works in the
ordinary course of nature and in the rain at Elijah's prayer or the water
made wine; ``when, for the admonition of men, they are thrust in by an
unusual changeableness, then they are called miracles'' (\emph{De Trin.}\
III.5.11--6.11). The magicians
of Pharaoh made serpents and failed at the third plague, the lice,
confessing ``This is the finger of God'' (Exod 8:19); so ``not even those
angels and powers of the air that transgressed \dots\ through whom magic
arts have whatever power they have, can do anything except by power given
from above'' (III.7.12). The angels are not creators, good or bad.
Hidden seeds of all things lie in the elements, \latin{occulta quaedam
semina in istis corporeis mundi huius elementis latent} (III.8.13), and
``as mothers are pregnant with young, so the world itself is pregnant with
the causes of things that are born'' (III.9.16);
the angels, knowing those seeds by the subtlety of their perception, may
apply them and hasten their growth, but
\begin{quote}
it is not right to think not only the bad but even the good angels to be
creators \dots\ But neither do the good angels do these things, except
as far as God commands, nor do the evil ones do them wrongfully, except as
far as He righteously permits. For the malignity of the wicked one makes
his own will wrongful; but the power to do so, he receives rightfully,
whether for his own punishment, or, in the case of others, for the
punishment of the wicked, or for the praise of the good. (III.8.13)
\end{quote}
The opening of the same chapter, that ``the matter of visible things is
subservient'' not ``to the bidding of those wicked angels; but rather to
that of God'' (III.8.13), is Aquinas's \emph{sed contra} for his article
that bodily matter does not obey the angels' mere will, and the doctrine
of the seeds is his \emph{sed contra} for the angels' power to move bodies
locally and his answer on the demons' wonders (\ST{I q.~110 aa.~2--3};
q.~114 a.~4 ad~2). His reply on the magicians follows the eighty-three
questions: \latin{magi per privatos contractus, boni christiani per
publicam iustitiam, mali christiani per signa publicae iustitiae} ---
magicians work wonders by private compacts, good Christians by public
justice, bad Christians by the signs of public justice (\emph{De div.\
qq.\ 83} q.~79.4; \ST{I q.~114 a.~4 ad~3}; \S\ref{sec:q114}).

On this ground Augustine determines that the appearances of God to the
patriarchs were wrought through angels. ``It is manifest, accordingly,
that all those appearances to the fathers, when God was presented to them
according to His own dispensation, suitable to the times, were wrought
through the creature. And if we cannot discern in what manner He wrought
them by ministry of angels, yet we say that they were wrought by angels''
(\emph{De Trin.}\ III.11.22). Scripture is his authority: the Epistle to
the Hebrews distinguishes the word ``spoken by angels'' from the salvation
``spoken by the Lord,'' and asks, ``Are they not all ministering spirits,
sent to minister for them who shall receive the inheritance of
salvation?'' (Heb 1:14; 2:2--3). Stephen tells that an angel appeared to
Moses in the bush (Acts 7:30), and the angel who called to Abraham on the
mountain swore ``By myself I have sworn, saith the Lord'' (Gen 22:16;
III.11.24--25).
Above all, the Law was received ``by the disposition of angels'' (Acts
7:53) and ``ordained by angels in the hand of a mediator'' (Gal 3:19); in
those angels ``were certainly the Father, and the Son, and the Holy
Spirit,'' figuratively signified by a creature and not seen in their
substance, and the Son ``was in some wonderful and unspeakable manner in
the angels, by whose disposition the Law itself was given'' (III.11.26).
Augustine names the contrary reading, of those who, because the Scripture
there specifies an angel, ``will have the Son of God Himself and in
Himself to be understood'' (III.11.23), and answers it with the plural of
the Epistle, ``by angels.'' The Catechism teaches that the angels
``communicated the law by their ministry'' (CCC~332). At Sinai ``great
marvels were wrought, by the ministry of angels, before the people on the
mount where the law was being given'' (\emph{De civ.\ Dei} X.13), and
God's word is heard
``by the mental ear of His ministers and messengers,'' and what they
receive ``they execute without delay or difficulty in the sensible and
visible world'' (\emph{De civ.\ Dei} X.15).

Augustine is as reserved about the manner of these works as he is sure of
the fact. Whether the angels, ``secretly working by the spiritual quality
of their body abiding still in them, assume somewhat from the inferior
and more bodily elements,'' which they change ``like a garment,'' or
change their own bodies, ``it reaches further than my purpose can carry
me to inquire'' (\emph{De Trin.}\ III.1.5; cf.\ III.10.21). He speaks at
times, with the learned of his day, of the angels' and demons' bodies as
of subtle or aerial substance, and he does not make it a point of faith:
if anyone holds that the devils have no bodies, ``this is not a matter
either to be laboriously investigated, or to be debated with keenness''
(\emph{De civ.\ Dei} XXI.10). Aquinas reads him so: Augustine ``often
makes use of this opinion in his books, although he does not mean to
assert it'' (\ST{I q.~54 a.~5}). The Church teaches that the angels are
``purely spiritual creatures'' (CCC~330), and Aquinas determines that they
have no bodies naturally united to them and assume bodies when they appear
(\ST{I q.~51 aa.~1--2}; \S\ref{sec:q51}); for the second article he
cites Augustine (\emph{De civ.\ Dei} XVI) that the angels appeared to
Abraham in assumed bodies.

Within the same government Augustine gives the demons' power and its
limits. The little book \emph{On the Divination of Demons} grew from a
conversation during the holy days of the octaves, when some claimed that the
demolition of the temple of Serapis at Alexandria had been foretold
(\emph{De divinatione daemonum} 1.1). Augustine answers that what God
permits is not thereby good (2.6), and then explains the demons' powers
from their nature:
\begin{quote}
\latin{Daemonum ea est natura, ut aerii corporis sensu terrenorum
corporum sensum facile praecedant; celeritate etiam propter eiusdem aerii
corporis superiorem mobilitatem non solum cursus quorumlibet hominum vel
ferarum, verum etiam volatus avium incomparabiliter vincant.} (3.7)
\end{quote}
Their aerial bodies give them keener senses and swifter movement than
any man, beast, or bird, and their long life a far greater experience;
yet no one prefers a dog to a good man for its scent (3.7). Their
foreknowledge is conjecture. They often foretell what they themselves
will do, since they receive power to send diseases, to corrupt the air,
and to persuade the perverse, \latin{per illam subtilitatem suorum
corporum corpora hominum non sentientium penetrando, seseque
cogitationibus eorum per quaedam imaginaria visa miscendo}, penetrating
men's bodies unperceived by that subtlety and mingling themselves with
their thoughts through certain imaginary visions, waking or sleeping
(5.9). They foretell from natural signs, as a physician foresees health
in the body, and they learn men's dispositions from signs the mind
expresses in the body (5.9). Of
that last claim Augustine later wrote, \latin{rem dixi occultissimam
audaciore asseveratione quam debui} --- I spoke of a most hidden thing
with bolder assertion than I ought. That such things reach the demons'
knowledge experience has shown; but whether by signs given from the body
or by some other, spiritual power can hardly or not at all be found out by
men
(\emph{Retractationes} II.30). Aquinas quotes both the treatise and the
retractation (\ST{I q.~57 a.~4}; \S\ref{sec:q57}). Their divination is
far from prophecy, \latin{quam Deus per sanctos Angelos suos et Prophetas
operatur}, which God works through his holy angels and prophets; when the
demons repeat what they have heard of God's disposition they neither
deceive nor are deceived, but in their own predictions they are often
both. They are deceived because an order from above overturns their
plans, and because what they foresee from natural causes is changed
\latin{ab Angelis Deo summo pie servientibus alia dispositione ignota
daemonibus}, by the angels who serve the most high God, by a disposition
unknown to the demons; and they deceive from envy, rejoicing in men's
errors (6.10).

The demons' present dwelling is the dark air. The commentary on Genesis
calls them \latin{aeria \dots\ animalia}, aerial living beings, and
supposes that if the transgressors with their prince, \latin{nunc
diabolo, tunc archangelo}, now the devil, then an archangel, were before
their sin in the upper and calmer air, it is no wonder that after sin
they were thrust down into this murk; for some of ours, \latin{nonnulli
nostri}, Augustine adds, do not think them to have been heavenly or
supercelestial angels (\emph{De Gen.\ ad litt.}\ III.10.14). If they bore
heavenly bodies before, these may have been changed in punishment into an
aerial quality, so that they can suffer from fire; and the dark air is to
them \latin{quasi carcer \dots\ usque ad tempus iudicii}, a kind of
prison until the time of judgment (III.10.15). Aquinas makes this the
\emph{sed contra} of his article on the demons' place of punishment
(\ST{I q.~64 a.~4}), and reads Augustine's ``upper atmosphere'' as said
of the angels who sinned only (q.~61 a.~4 ad~2).

\subsubsection{The worship of the one God}\label{sec:aug-worship}

The eighth to tenth books of the \emph{City of God} are Augustine's
treatise on the cult of spirits. The Platonists, whom he counts the
noblest of the philosophers, divided rational souls into gods in heaven,
demons in the air, and men on earth, and held with Apuleius that the
demons carry the prayers of men to the gods and the gifts of the gods to
men, since ``no god has intercourse with man'' (VIII.14, 18, 20).
Apuleius defined them: ``Demons are of an animal nature, passive in soul,
rational in mind, aerial in body, eternal in time'' (VIII.16). Augustine
turns the definition against them: better bodies are no title to worship
(VIII.15), and passion of soul is the mark of misery, so that they are
``aerial animals who
are only rational that they may be capable of misery, passive that they
may be actually miserable, and eternal that it may be impossible for them
to end their misery!'' (VIII.16). The demons, then, are not messengers
between gods and men:
\begin{quote}
we must believe them to be spirits most eager to inflict harm, utterly
alien from righteousness, swollen with pride, pale with envy, subtle in
deceit; who dwell indeed in this air as in a prison, in keeping with
their own character, because, cast down from the height of the higher
heaven, they have been condemned to dwell in this element as the just
reward of irretrievable transgression. (VIII.22)
\end{quote}
What we may share with the holy angels is not their mediation but their
will. We come to their benevolence ``through resembling them in the
possession of a good will, through which we are with them, and live with
them, and worship with them the same God \dots\ it is not in locality we
are distant from them, but in merit of life'' (VIII.25). The Church's honour of the martyrs
follows the same rule: no priest at an altar over a martyr's body says
``I offer to thee a sacrifice, O Peter, or O Paul, or O Cyprian,'' for
sacrifice is offered there to ``the God who made them both men and
martyrs, and associated them with holy angels in celestial honor''
(VIII.27).

The ninth book answers those who distinguish good demons from bad. The
holy angels feel no anger when they punish, no fellow-suffering when they
relieve, no fear when they aid; Scripture ascribes such emotions to them
``because, though they have none of our weakness, their acts resemble the
actions to which these emotions move us'' (IX.5); Aquinas, citing the
book, says that none of these is said of the angels by way of passion
(\ST{I q.~59 a.~4 ad~2}). A mediator
between miserable mortals and the blessed immortals would have to share
something with each; the good angels cannot be such mediators, and the
evil angels can be, but only to separate:
\begin{quote}
Good angels, therefore, cannot mediate between miserable mortals and
blessed immortals, for they themselves also are both blessed and
immortal; but evil angels can mediate, because they are immortal like the
one party, miserable like the other. To these is opposed the good
Mediator, who, in opposition to their immortality and misery, has chosen
to be mortal for a time, and has been able to continue blessed in
eternity. \dots\ There is, then, a wicked mediator, who separates
friends, and a good Mediator, who reconciles enemies. (IX.15)
\end{quote}
Christ is Mediator as man, not as Word, and he does not lead us to the
angels to be blessed by sharing their nature, ``but He leads us straight
to that Trinity, by participating in which the angels themselves are
blessed. Therefore, when He chose to be in the form of a servant, and
lower than the angels, that He might be our Mediator, He remained higher
than the angels, in the form of God'' (IX.15; 1 Tim 2:5). The Church
confesses Christ ``the sole mediator and way of salvation'' (Paul~VI,
profession of faith \emph{Solemni hac liturgia}, the Credo of the People
of God, 1968, n.~23). Scripture itself refuses the name
of demon to the good: ``never have we read in Scripture of good demons''
(IX.19). The name, taken from knowledge, fits them, for ``Knowledge
puffeth up: but charity edifieth'' (1 Cor 8:1): ``The demons, then, have
knowledge without charity'' (IX.20). When they cried, ``What have we to do
with thee, Jesus of Nazareth? Art thou come to destroy us?'' (Mark 1:24),
``it is clear that they had great knowledge, and no charity''; Christ made
himself known to them ``not as to the holy angels, who know Him as the
Word of God, and rejoice in His eternity, which they partake, but as was
requisite to strike with terror,'' by temporal effects of his power
(IX.21). Aquinas cites this chapter for the measure of the demons'
knowledge of the Incarnation (\ST{I q.~64 a.~1 ad~4}). The Platonists may
call the holy angels gods, since Scripture calls even men gods (Ps
81[82]:6); but ``the name of demon is so detestable that we cannot bear in
any sense to apply it to the holy angels'' (IX.23).

The tenth book shows that the holy angels themselves desire worship to be
given to God alone. The worship called in Greek \gk{latreia}, ``the
service due to God only,'' belongs ``only to that God who is the true God,
and who makes His worshippers gods'' (X.1). ``If he does not
worship God, he is wretched, because deprived of God; if he worships God,
he cannot wish to be worshipped in God's stead'' (X.3). The true
sacrifice is every work done that we may cleave to God in holy
fellowship, and ``the whole redeemed city \dots\ is offered to God as our
sacrifice through the great High Priest'' in the sacrament of the altar
(X.6). Of that sacrifice the angels are part:
\begin{quote}
It is very right that these blessed and immortal spirits, who inhabit
celestial dwellings, and rejoice in the communications of their Creator's
fullness, firm in His eternity, assured in His truth, holy by His grace,
since they compassionately and tenderly regard us miserable mortals, and
wish us to become immortal and happy, do not desire us to sacrifice to
themselves, but to Him whose sacrifice they know themselves to be in
common with us. For we and they together are the one city of God, to which
it is said in the psalm, ``Glorious things are spoken of thee, O city of
God;'' the human part sojourning here below, the angelic aiding from
above. (X.7)
\end{quote}
In the Latin they know themselves to be God's sacrifice \latin{nobiscum},
with us, and the heavenly city is a council-house where \latin{geritur
namque ibi cura de nobis}, care is taken for us. When the angels hear our
prayers, ``it is He Himself who hears us in them, as in His true temple
not made with hands'' (X.12). Between angels who demand sacrifice and
angels who direct it to God the choice is plain, for those who direct it
to God ``prove how sincerely they love us, since they wish by sacrifice to
subject us, not to themselves, but to Him by the contemplation of whom
they themselves are blessed'' (X.16). When we offer
ourselves to God in sacrifice,
\begin{quote}
It is then that the angels, and all those superior powers who are mighty
by their goodness and piety, regard us with pleasure, and rejoice with us
and assist us to the utmost of their power. But if we offer such worship
to them, they decline it; and when on any mission to men they become
visible to the senses, they positively forbid it. (X.19)
\end{quote}
Scripture gives the examples: the angel of the Apocalypse says, ``See thou
do it not. I am thy fellow servant \dots\ Adore God'' (Apoc 19:10;
22:9), and Paul and Barnabas, imitating the holy angels, refused the
sacrifice of the Lycaonians (Acts 14). The proud
spirits who exact worship delight not in the smoke of victims but ``with
the suppliant spirit which they deceive and hold in subjection'' (X.19).
Christ, who in the form of God receives sacrifice with the Father, ``in
the form of a servant He chose rather to be than to receive a sacrifice,''
and ``He designed that there should be a daily sign of this in the
sacrifice of the Church, which, being His body, learns to offer herself
through Him'' (X.20). The power granted to the demons to persecute the
city is ``not merely harmless, but even useful to the Church, completing
as it does the number of martyrs'' (X.21), and men of God cast out the
power of the air ``by exorcising it, not by propitiating it'' (X.22).

The short treatise \emph{On True Religion} puts the rule of the angels'
honour in a sentence. Whatever the highest angel worships, the lowest man
must worship, \latin{Quod ergo colit summus angelus, id colendum est
etiam ab homine ultimo}, for angel and man are wise and truthful from the
one unchangeable Wisdom and Truth; the best angels wish us to worship
with them the one God by whose contemplation they are blessed. We love
them, and do not envy them their readier joy, but love them the more
because we are bidden to hope for the same:
\begin{quote}
\latin{Quare honoramus eos caritate, non servitute. Nec eis templa
constituimus: nolunt enim se sic honorari a nobis; quia nos ipsos cum boni
sumus, templa summi Dei esse noverunt.} (\emph{De vera religione} 55.110)
\end{quote}
We honour them with love, not with servitude; we build them no temples,
for they do not wish to be so honoured, knowing that we ourselves, when we
are good, are temples of the most high God; and rightly the angel forbade
the man to adore him, being his fellow servant. Augustine adds the
confidence of the communion of saints: \latin{Quisquis Angelorum diligit
hunc Deum, certus sum quod etiam me diligit. Quisquis in illo manet, et
potest humanas preces sentire, in illo me exaudit} --- whatever angel
loves this God, I am sure that he loves me too; whoever abides in him and
can perceive human prayers hears me in him (55.112). The Church joins
with the angels to adore the thrice-holy God and invokes their assistance
(CCC~335); she venerates the heavenly spirits and ``has recourse to their
prompt intercession'' (Directory on Popular Piety and the Liturgy 215;
\S\S\ref{sec:liturgy}, \ref{sec:devotion}).

\subsubsection{The sons of God}\label{sec:aug-gen6}

``The sons of God seeing the daughters of men, that they were fair, took
to themselves wives of all which they chose'' (Gen 6:2, Douay). Augustine
reads the verse otherwise than many before him (\S\ref{sec:fathers-before}).
He knows their reading, ``whence many suppose that they were not men but
angels,'' since in some copies the same persons are called angels of God
(\emph{De civ.\ Dei} XV.22), and he refuses it. That angels appeared in
bodies that could be seen and touched Scripture testifies; whether some
spirits embodied in air are capable of lust, as many reports affirm, ``I
dare not determine''; ``but certainly I could by no means believe that
God's holy angels could at that time have so fallen, nor can I think that
it is of them the Apostle Peter said, `For if God spared not the angels
that sinned, but cast them down to hell' \dots\ I think he rather speaks
of these who first apostatized from God, along with their chief the
devil, who enviously deceived the first man under the form of a serpent''
(XV.23). Scripture calls godly men angels, as John the Baptist (Mark 1:2)
and the priest (Mal 2:7). Giants were on the earth before the marriages
as well as after (Gen 6:4). And the text itself decides: after the
angels of God took wives, ``the Lord God said, My Spirit shall not always
strive with these men, for that they also are flesh'':
\begin{quote}
For by the Spirit of God they had been made angels of God, and sons of
God; but declining towards lower things, they are called men, a name of
nature, not of grace; and they are called flesh, as deserters of the
Spirit, and by their desertion deserted [by Him]. (XV.23)
\end{quote}
The formula answers the one in the sermon on Psalm~103: there ``angel''
names an office, here ``men'' names a nature and ``sons of God'' a grace.
The writings under the name of Enoch that tell of angelic fathers of the
giants he sets aside: Enoch left ``some divine writings,'' as Jude
attests (Jude 14), but the books produced under his name ``are properly
judged by prudent men to be not genuine'' and stand outside the canon
(XV.23). His conclusion is that ``the sons of
God, who were according to the flesh the sons of Seth, sunk into this
community when they forsook righteousness'' (XV.23). Aquinas quotes the
chapter for this conclusion (\ST{I q.~51 a.~3 ad~6}; \S\ref{sec:q51}).

\subsubsection{The angels in the heavenly city}\label{sec:aug-vision}

The last book of the \emph{City of God} returns to the angels at its
beginning and near its end. In the eternal city ``all the citizens shall
be immortal, men now for the first time enjoying what the holy angels
have never lost'' (XXII.1). God made nothing better than the spirits to
whom he gave intelligence and made ``capable of contemplating and
enjoying Him, and united in our society, which we call the holy and
heavenly city, and in which the material of their sustenance and
blessedness is God Himself, as it were their common food and
nourishment.'' Those he foreknew would fall he left free, ``deeming it
to be more befitting His power and goodness to bring good out of evil than
to prevent the evil from coming into existence''; and he ``rewarded those
who continued in their attachment to the
supreme good with the assurance of endless stability as the meed of their
fidelity'' (XXII.1).

The angels already see what we hope to see. The peace of God passes all
human understanding, and the angels understand it in their own measure;
of the face-to-face vision promised to us (1 Cor 13:12) Augustine says:
\begin{quote}
Such also is now the vision of the holy angels, who are also called our
angels, because we, being rescued out of the power of darkness, and
receiving the earnest of the Spirit, are translated into the kingdom of
Christ, and already begin to belong to those angels with whom we shall
enjoy that holy and most delightful city of God of which we have now
written so much. Thus, then, the angels of God are our angels, as Christ
is God's and also ours. They are God's, because they have not abandoned
Him; they are ours, because we are their fellow-citizens. (XXII.29)
\end{quote}
The saying of the Lord, ``their angels in heaven always see the face of my
Father'' (Matt 18:10), is read here of that fellowship: ``As, then, they
see, so shall we also see; but not yet do we thus see'' (XXII.29). The
\emph{Enchiridion} says the same of the peace of heaven. It passes our
understanding, ``not that of those who always behold the face of their
Father''; ``But when we shall be equal unto the angels of God then we
shall see face to face, as they do; and we shall have as great peace
towards them as they have towards us, because we shall love them as much
as we are loved by them'' (\emph{Enchiridion} 63). In that city there
will be degrees of honour and glory, and no envy, ``as now the archangels
are not envied by the angels'' (\emph{De civ.\ Dei} XXII.30). Aquinas
determines that men are taken up into the angelic orders by grace, not by
nature, on the promise that the children of the resurrection are equal to
the angels, and rejects a separate society of the saved on Augustine's
word that there is one society of men and angels (\ST{I q.~108 a.~8};
\S\ref{sec:q108}).

\subsubsection{The Augustinian loci and their reception}\label{sec:aug-loci}

\begin{fourstable}{Work and locus}{Teaching}{Received by Aquinas at}{Grade}
\emph{Enarr.\ in Ps.} 103, s.~1.15; \emph{De civ.\ Dei} XV.23; \emph{De Gen.\ ad litt.}\ V.19.37 & ``Angel'' names the office of messenger; by nature they are spirits; the whole heavenly city bears the name & Parallel: \ST{I q.~108 a.~5 ad~1} (\latin{Angelus nuntius dicitur}) & Church teaching (CCC~329 quotes the \emph{Enarratio})\\
\emph{De civ.\ Dei} XI.9, 11, 32--34; \emph{De Gen.\ ad litt.}\ II.8.16; IV.22.39; VIII.20.39 & The angels are made by God, signified by the light of the first day, formed by turning to the Word, created with the world & \ST{I q.~61 a.~1 ad~1} (XI.9); q.~61 a.~2 ad~2 (VIII.20); q.~61 a.~3; q.~62 a.~3 obj.~1 and ad~1 (II.8) & Defined that they are created from nothing with the corporeal creature (DS~800); Disputed whether before the visible world (q.~61 a.~3)\\
\emph{De Gen.\ ad litt.}\ II.8.16--19; IV.22.39--35.56; \emph{De civ.\ Dei} XI.7, 29 & Knowledge of things in the Word (day, morning) and in their own nature (evening), referred always to praise, held at once & \ST{I q.~55 a.~2 ad~1}; q.~56 a.~2; q.~58 a.~1 s.c.; q.~58 a.~2 s.c.; q.~58 aa.~6--7 s.c. & Common teaching (\S\ref{sec:q58})\\
\emph{De Gen.\ ad litt.}\ V.19.38--39; \emph{De civ.\ Dei} IX.20--22 & The angels knew the mystery of the kingdom from the ages; the demons know by conjecture and by temporal effects, and are often deceived & \ST{I q.~57 a.~5 obj.~1}; q.~64 a.~1 obj.~4, ad~4 & Common teaching\\
\emph{De civ.\ Dei} XI.16, 19--20, 33; XII.1; \emph{De Gen.\ ad litt.}\ XI.15.20 & Two societies of angels, alike good by nature, divided by will; with men, two cities and not four & \ST{I q.~63 a.~5 ad~2} (XI.19); q.~64 a.~1 obj.~5 (XI.33); q.~108 a.~1 (XII.1) & Defined that the demons were created good and became evil by their own will (DS~800); Common teaching on the two cities\\
\emph{De civ.\ Dei} XII.6--8 & The evil will has no efficient cause but a deficient one, the defection of the will from the highest good & Parallel: \ST{I q.~63 a.~1 ad~4} (a good sought without its rule) & Common teaching\\
\emph{De civ.\ Dei} XII.9 & The good will created with the nature; grace given with it; the fallen had less grace, or the others more help; charity poured out by the Holy Spirit & \ST{I q.~62 a.~3 s.c.}; q.~60 a.~5 obj.~4; q.~108 a.~4 obj.~2; q.~108 a.~8 & Thomist position that the angels were created in sanctifying grace (q.~62 a.~3); the Lombard's school dissents (\S\ref{sec:q62})\\
\emph{De civ.\ Dei} XI.13--15, 17; XIV.3 & The devil did not stand in the truth, was not made evil, sinned from the beginning of his sin; proud and envious, not carnal & \ST{I q.~63 a.~2 s.c.} (XIV.3); a.~5 co., ad~1, ad~2, ad~4; a.~6 s.c. (XI.15) & Defined that the devil was created good and became evil of himself (DS~800)\\
\emph{De Gen.\ ad litt.}\ XI.13.17--27.34 & Pride before envy; the devil fell from what he would have received; possibly at the very beginning; four possibilities; he tempts only as permitted & Order of pride and envy: \ST{I q.~63 a.~2}; the moment determined otherwise: q.~63 aa.~5--6; the lower-order opinion: q.~63 a.~7 & Common teaching that pride was the first sin; Disputed whether the devil ever stood in grace (Appendix~\ref{app:first-sin})\\
\emph{Enchir.}\ 28--29, 61--62; \emph{De civ.\ Dei} XI.13; XXII.1 & The fall irrevocable and the rest confirmed; the ruin repaired from redeemed men; the number known to God & \ST{I q.~63 a.~9 ad~3}; q.~108 a.~8 & Church teaching on irrevocability (CCC~393); Common teaching on the repair of the ruin\\
\emph{Enchir.}\ 56--60; \emph{De civ.\ Dei} XXII.30 & The Church includes the holy angels; the orders named, their differences unknown; archangels above angels & Ordering given after Dionysius and Gregory: \ST{I q.~108 aa.~5--6} & Common teaching that the angels belong to the one Church; Disputed on the ordering (\S\ref{sec:q108})\\
\emph{De Trin.}\ III.4.9--9.18; \emph{De Gen.\ ad litt.}\ VIII.24.45--25.47; \emph{De div.\ qq.\ 83} q.~79 & The bodily creation governed through the spiritual; the angels not creators; the seeds of things; the demons' power only by permission & \ST{I q.~110 a.~1 s.c.}\ and co.; a.~2 s.c.; a.~3 s.c.; q.~108 a.~6; q.~109 a.~4 s.c.; q.~114 a.~4 ad~2--3 & Common teaching\\
\emph{De Trin.}\ III.10.19--11.27; \emph{De civ.\ Dei} X.13, 15 & The theophanies and the Law wrought by angels, God speaking by them & --- & Church teaching that the Law was given by the angels' ministry (CCC~332; Acts 7:53; Gal 3:19); Disputed whether the Son appeared in his own person (De Trin.\ III.11.23)\\
\emph{De div.\ daem.}\ 3.7--6.10; \emph{Retract.}\ II.30; \emph{De Gen.\ ad litt.}\ III.10.14--15; \emph{De civ.\ Dei} XXI.10 & The demons' keen sense, speed, and experience; conjectural foreknowledge; thoughts known by signs, with reserve; the dark air their prison & \ST{I q.~57 a.~4}; q.~54 a.~5; q.~61 a.~4 ad~2; q.~64 a.~4 s.c. & Common teaching on conjectural knowledge; Church teaching that the angels are purely spiritual (CCC~330); Thomist position that they have no bodies (q.~51 a.~1)\\
\emph{De civ.\ Dei} VIII.14--IX.23 & The demons false mediators; Christ the one Mediator; the good angels not called demons & \ST{I q.~59 a.~4 ad~2} (IX.5); q.~64 a.~1 ad~4 (IX.21); q.~63 a.~4 ad~1 (X.11) & Church teaching (Paul~VI, \emph{Credo} 23; 1 Tim 2:5)\\
\emph{De civ.\ Dei} X.1--7, 12, 16, 19--22; \emph{De vera rel.}\ 55.110--112 & Sacrifice to God alone; the angels offered with us and refusing worship; honour of love, not servitude & --- & Church teaching that the Church adores God with the angels and invokes them (CCC~335; Directory 215); Common teaching on love, not servitude\\
\emph{De civ.\ Dei} XV.22--23 & The sons of God are the sons of Seth, not angels & \ST{I q.~51 a.~3 ad~6} & Common teaching in the West after Augustine; the earlier Fathers dissent (\S\ref{sec:fathers-before})\\
\emph{De civ.\ Dei} XXII.1, 29--30; \emph{Enchir.}\ 63 & The angels see God face to face; they are our angels as fellow-citizens; men equal to them, in degrees without envy & \ST{I q.~108 a.~8} & Church teaching that the angels behold the face of the Father (CCC~329; Matt 18:10); Common teaching on degrees of glory\\
\end{fourstable}
